% !TeX root = théorie.tex

\section{Groupes abéliens}

\begin{definition}[Jonction par le puits]
Soient $\left(G_i\right)_{1 \leqslant i \leqslant n} =
\left(S_i, A_i, p_i\right)_{1 \leqslant i \leqslant n}$ des graphes des tas de
sable où $\left(S_i\right)_{1 \leqslant i \leqslant n}$ sont des ensembles
disjoints.
La jonction par le puits est le graphe des tas de sable
$(\left\{p\right\} \cup \bigcup_{i=1}^{n} \widetilde{S_i},
\bigcup_{i=1}^{n} A_i, p)$ où tous les $\left(p_i\right)_
{1 \leqslant i \leqslant n}$ sont identifiés au puits commun $p$.
\end{definition}

\begin{proposition}
Si $G$ est la jonction par le puits de $\left(G_i\right)
_{1 \leqslant i \leqslant n}$ alors $\mathcal{M}\left(G\right)
= \mathcal{M}\left(G_1\right) \times \ldots \times \mathcal{M}\left(G_n\right)$
et $\mathcal{G}\left(G\right)
= \mathcal{G}\left(G_1\right) \times \ldots \times \mathcal{G}\left(G_n\right)$.
\end{proposition}

\begin{proof}
À faire
\end{proof}

\begin{theoreme}
Soit $(G, \cdot)$ un magma fini.
On a l'équivalence entre :
\begin{itemize}[label=\textbullet]
    \item $(G, \cdot)$ est isomorphe à un groupe des tas de sable récurrents
    \item $(G, \cdot)$ est un groupe abélien
\end{itemize}
\end{theoreme}

\begin{proof}
Le sens direct découle de la construction du groupe des tas de sable récurrents.
Supposons que $(G, \cdot)$ est un groupe abélien.
D'après le théorème de Kronecker, il existe des entiers naturels non nuls
$a_1 | a_2 | \ldots | a_k$ tels que :
$$
G \simeq \mathbb {Z} / a_1 \mathbb{Z} \times \mathbb{Z} / a_2 \mathbb{Z} \times \cdots \times \mathbb{Z} / a_k \mathbb{Z}
$$
Pour tout $i \in [\![1;k]\!]$, posons le graphe des tas de sable $H_i =
\left( \left\{s_i, p_i\right\}, \left(\left\{\left(s_i, p_i\right)\right\},
\left(s_i, p_i\right) \mapsto a_i\right), p_i \right)$
de sorte que $\mathcal{M}\left(H_i\right) = \mathcal{G}\left(H_i\right) =
\mathbb{Z} / a_i \mathbb{Z}$.
Posons $H$ la jonction par le puits de $\left(H_i\right)_
{1 \leqslant i \leqslant n}$ de sorte que $\mathcal{G}(H) \simeq G$ d'après la
proposition ci-dessus.
\end{proof}

\begin{definition}[Pseudo-base]
Soit $(M, \cdot, e)$ un monoïde.
Soit $F \subseteq M$ et $E \in \mathcal{MP}(F)$.
$(F, M)$ est une pseudo-base du monoïde  si :
\begin{enumerate}
\item $\forall x \in M, \exists ! E'_x \subseteq E: x = \prod_{f \in E'_x} f$
\item $\forall (x, y) \in M^2, |E'_{xy}| \leqslant |E'_x| + |E'_y|$
\item $\forall f \in F, \exists k \in \mathbb{N}^*: |E'_{f^k}| < k$
\end{enumerate}
avec la convention $\prod_{y \in \emptyset} y = e$.
$F$ est alors une partie génératrice.
\end{definition}

\begin{proposition}
Soit $(M, *)$ un magma fini.
on a l'équivalence entre :
\begin{itemize}[label=\textbullet]
    \item $(M, *)$ est isomorphe à un monoïde des tas de sable stables
    \item $(M, *)$ est un monoïde abélien admettant une pseudo-base
\end{itemize}
\end{proposition}

\begin{proof}
Le sens direct découle de la construction du groupe des tas de sable stables.
Réciproquement, soit $(M, *)$ un monoïde abélien admettant une pseudo-base
$(F, E)$.
Construisons un multi-graphe des tas de sable $G = (S, A, p)$.
Soit $p$ un puits. Posons $S = F \cup \{p\}$.
Pour tout élément $f \in F$, par définition d'une pseudo-base, il existe
$\hat{E}_f \subseteq E : f^{m_E(f)+1} = \prod_{g \in \hat{E}_f} g$ (d'après la
condition 1) et $|\hat{E}_f| \leqslant m_E(f) + 1$ (d'après la condition 2).
Posons $A = \left\{ (f, g, m) \;|\; f \in F, (m, g) \in \hat{E}_f \right\}
\cup \left\{ (f, p, m_E(f) + 1 - |\hat{E}_f|) \;|\; f \in F \right\}$.
La condition $3$ garantit que le puits est accessible.

L'application $\varphi: \begin{matrix}
M \rightarrow \mathcal{M}(G) \\
x = \prod_{f \in E'_x} f \mapsto \sum_{f \in E'_x} \mathbf{1}_f
\end{matrix}$
est un morphisme de monoïdes.
\end{proof}

\begin{corollaire}
Tout monoïde abélien admettant une pseudo-base est simplifiable.
\end{corollaire}

%\begin{question}
%Soit $(M, *)$ un magma fini.
%A-t-on l'équivalence entre :
%\begin{itemize}[label=\textbullet]
%    \item $(M, *)$ est isomorphe à un monoïde des tas de sable stables
%    \item $(M, *)$ est un monoïde abélien simplifiable
%\end{itemize}
%ou faut-il ajouter des conditions pour le second point ?
%Je donne ci-dessous un début de démonstration inachevé.
%\end{question}

%\begin{proof}
%Le sens direct découle de la construction du groupe des tas de sable stables.
%Réciproquement, soit $(M, *)$ un monoïde abélien.
%L'ensemble des parties génératrices de $(M, *)$ est non vide car $M$ en est une.
%Il existe alors $X$ une partie génératrice minimale au sens de l'inclusion.
%Notons $X = \left\{ x_1, x_2, \ldots, x_n \right\}$.
%Posons $S = \left\{ s_1, s_2, \ldots, s_n, p \right\}$.

%Munissons $\mathcal{MP}(X)$ de l'ordre total :
%$$
%(X, m) \sqsubseteq (X, m')
%\iff \left(|(X, m')|, m(x_1), m(x_2), \ldots, m(x_n)\right)
%\leqslant \left(|(X, m')|, m'(x_1), \ldots, m'(x_n)\right)
%$$
%Pour tout $i \in [\![1;n]\!]$, posons
%$(K_i, \xi_i) = \min \left\{ (k, \chi) \in \mathbb{N}^* \times
%\mathcal{MP}(X) \;|\; \begin{matrix}
%x_i^k = \prod_{x \in \chi} x \\
%\chi \neq \left\{\left(x_i, k\right)\right\} \\
%\end{matrix} \right\}$.
%Posons $A = \left(S, \begin{matrix}
%\left(s_i, s_j\right) \mapsto m_{\xi_i}(x_j) \\
%\left(s_i, p \right) \mapsto K_i - \left|\xi_i\right|
%\end{matrix} \right)$.
%Posons $G = (S, A, p)$
%Montrons maintenant que $\mathcal{M}(G) \simeq M$.

%Comme $(M, *)$ est un semi-groupe abélien, tout élément $y \in M$ peut s'écrire
%sous la forme $y = \prod_{x \in \chi} x$ avec $\chi \in
%\mathcal{MP}(X)$.
%Posons alors $\varphi(y) = \sum_{i=1}^{n} m_\chi(x_i) . \mathbf{1}_{s_i}$.
%Montrons que l'application $\varphi$ est bien définie c'est-à-dire que la valeur
%de $\varphi(y)$ ne dépend pas du choix de $\chi$.
%\end{proof}

\begin{remarque}
Il est possible que la démonstration pour les groupes abéliens soit un cas particulier du problème précédent.
\end{remarque}

\begin{definition}
Soit $(G, \cdot)$ un groupe.
Posons
$$
a(G) = \min \left\{ |A| \;|\; \mathcal{G}(S, A, p) \simeq G \right\}
$$
\end{definition}

\begin{remarque}
D'après la preuve précédente, $a(G) \leqslant \sum_{i = 1}^k  a_i$.
\end{remarque}

\begin{question}
$$
\normalfont{\textsc{Problème sans nom}} :
\left\{ \begin{array}{lcl}
\text{\textbf{Entrée}} & \textbf{ : } & \text{un groupe abélien } (G, \cdot) \\
\text{\textbf{Sortie}} & \textbf{ : } & a(G)
\end{array} \right.
$$
Dans le problème ci-dessus, $(G, \cdot)$ est représenté par le multi-ensemble de
sa décomposition en groupes cycliques.

Ce problème est \textbf{NP}. Est-il \textbf{NP}-difficile ?
\end{question}